% Unidad U012. Biografía estructural completa de pi.

\label{gfp:b:chap:biografia-pi-autonoma}

\section{Causal separation between seed, character, and evaluation}

The construction is divided into three strata:
\[
 \text{finite selection}
 \longrightarrow
 \text{construction of the closure character}
 \longrightarrow
 \text{Archimedean evaluation}.
\]
The decimal expansion does not intervene in the first stratum. The type chain
is
\begin{equation}
 104\,976
 \longrightarrow468\,\mathcal U_6
 \xrightarrow{\Psi}243\,\mathcal W_6
 \lhook\joinrel\longrightarrow\mathbb F_3^6,
 \qquad|\mathbb F_3^6|=729.
 \label{gfp:b:eq:cadena-tipica-pi-autonoma}
\end{equation}
The cardinalities count, respectively, enriched initial conditions,
distinct emissions, visible ternary words, and elements
of the complete ambient fiber.

For \(U_6=(u_1,\ldots,u_6)\), the projection is
\[
 \Psi(U_6)=((-u_1)\bmod3,\ldots,(-u_6)\bmod3).
\]
The census of the \(43\) orbits of \(D_3\) proves that the closure
orbit is the only one whose six orientations possess at once:
\[
 \text{five lifts},\qquad
 \text{mass }1008,\qquad
 1008=432+4\cdot144.
\]
The gauge
\[
 \operatorname{diag}_c=6,
 \qquad
 \operatorname{card}=ES
\]
select a single representative:
\begin{equation}
 \boxed{w_6^\pi=010211.}
 \label{gfp:b:eq:semilla-pi-autonoma}
\end{equation}
The selection usis catalogue, multiplicitiis, orbital action, and orientation;
it dois not consult to prefix of \(\pi\).

\section{Lifts \texorpdfstring{\(L_0,L_1\)}{L0, L1} and \(1^{\mathrm a}\) coinductive boundary}

On vectors row of \(\mathbb F_3^6\), let
\[
L_0=
\begin{pmatrix}
2&2&2&1&2&1\\
2&1&2&2&1&1\\
1&0&0&1&0&2\\
0&0&1&1&1&0\\
2&2&2&0&2&1\\
2&1&2&1&0&0
\end{pmatrix},
\quad
L_1=
\begin{pmatrix}
2&1&2&0&2&2\\
1&2&1&1&2&0\\
1&2&1&1&1&2\\
0&1&0&0&2&1\\
2&0&2&1&0&1\\
0&2&2&2&1&1
\end{pmatrix}.
\]
If \(b_1=w_6^\pi\), the recurrence of blocks is
\[
 b_{k+1}=b_kL_{\varepsilon_k},
 \qquad
 (\varepsilon_1,\varepsilon_2,\varepsilon_3,\varepsilon_4)=(0,0,1,1),
\]
and the accumulated word is formed by concatenation:
\[
 w_{6(k+1)}=w_{6k}\mathbin{\|}b_{k+1}.
\]
A word of length \(6k\) is not multiplied by a matrix \(6\times6\).

For the closing channel:
\[
\begin{array}{c|c}
\text{depth}&\text{new block}\\ \hline
6&010211\\
12&012222\\
18&010211\\
24&002111\\
30&110221
\end{array}
\]
and
\begin{equation}
 \boxed{
 w_{30}^{\pi}
 =010211|012222|010211|002111|110221.}
 \label{gfp:b:eq:w30-pi-autonomo}
\end{equation}
The following boundary is
\[
 u_5^\pi=222220.
\]
Along with the other channels,
\[
 R_{36}=
 \begin{pmatrix}
 222220\\021222\\102011
 \end{pmatrix}.
\]
The symbol \(R_{36}\) records the sixth block of three histories and opens a
boundary; it is not a terminal.

The certified branch of twenty blocks begins
\begin{align*}
\pi:\;&010211|012222|010211|002111|110221|222220|111201|\\
&212121|200121|100100|101222|022212|012012|111210|\\
&121011|200220|120210|000101|022010|020111.
\end{align*}
For this reason, ni \(w_{30}\) ni \(R_{36}\) terminate the genealogy.

\section{Closure microfiber of \(\pi\): \(5\) oriented regions and multiplicity \(432+4\cdot144=1008\)}

The microfiber
\[
 \mathscr R_\pi^{\rm micro}
 :=\{U\in\mathcal U_6:\Psi(U)=010211\}
\]
contains exactly five regions:
\[
\begin{array}{c|c|c|c}
U_6&E_{\rm sig}&\operatorname{mult}&\text{role}\\ \hline
501|614|498|169|272|272&555555&432&\perp\\
810|923|870|169|272|272&339555&144&++\\
810|983|810|169|272|272&393555&144&++\\
870|923|810|169|272|272&933555&144&++\\
870|923|810|418|674|521&933717&144&--
\end{array}
\]
Consequently,
\begin{equation}
 \boxed{
 5=1_\perp+3_{++}+1_{--},
 \qquad
 1008=432+4\cdot144.}
 \label{gfp:b:eq:descomposiciones-cinco-regiones-pi}
\end{equation}
The first identity classifies orientation functions. The second counts
provenance multiplicities. They do not count five indistinguishable copies.

The transversal region is
\[
 501|614|498|169|272|272,
 \qquad E_{\rm sig}=555555,
 \qquad\operatorname{mult}=432.
\]
The mutated return is
\[
 870|923|810|418|674|521,
 \qquad E_{\rm sig}=933717,
 \qquad\operatorname{mult}=144.
\]
Interchanging them would preserve the census \(3/1/1\) but falsify the complete
regional structure.

\section{Phase transfer on the multigraph of complete blocks}

For \(U_6=(U_1,\ldots,U_6)\), let
\[
 A(U_6)=(U_1,U_2,U_3),
 \qquad
 B(U_6)=(U_4,U_5,U_6).
\]
Each emission defines an edge between \(A(U_6)\) and \(B(U_6)\). The raw
multigraph has \(96\) vertices, distributed into components of sizes
\[
 28,28,28,4,4,4.
\]
The internal phase is quotiented by
\[
 (a,b,c)\sim(b,c,a)\sim(c,a,b).
\]
The representative is the lexicographically minimal rotation. The quotient
graph has \(32\) vertices and two components of sizes \(28\) and \(4\).

The anchors are
\[
\begin{aligned}
U_{\rm APP}&=903|850|850|252|109|252,\\
U_e&=850|963|890|149|252|212,\\
U_\varphi&=830|943|830|109|252|252.
\end{aligned}
\]
\Needspace{9\baselineskip}
We number the regions
\[
\begin{aligned}
U_\pi^{(1)}&=810|923|870|169|272|272,\\
U_\pi^{(2)}&=810|983|810|169|272|272,\\
U_\pi^{(3)}&=870|923|810|169|272|272,\\
U_\pi^{(4)}&=870|923|810|418|674|521,\\
U_\pi^{(5)}&=501|614|498|169|272|272.
\end{aligned}
\]

\section{\(5\) oriented closure cycles of the microfiber}

\subsection{Closure of \texorpdfstring{\(U_\pi^{(1)}\)}{U-pi 1}}
\begin{align*}
&252|109|252|903|850|850\to
903|850|850|252|109|252\to\\
&252|109|252|923|870|810\to
810|923|870|169|272|272\to\\
&272|169|272|963|890|850\to
850|963|890|149|252|212\to\\
&212|149|252|943|830|830\to
830|943|830|109|252|252\to
252|109|252|903|850|850.
\end{align*}

\Needspace{10\baselineskip}
\subsection{Closure of \texorpdfstring{\(U_\pi^{(2)}\)}{U-pi 2}}
\begin{align*}
&903|850|850|252|109|252\to
272|169|272|903|850|850\to\\
&272|169|272|983|810|810\to
810|983|810|169|272|272\to\\
&272|169|272|963|890|850\to
850|963|890|149|252|212\to\\
&212|149|252|943|830|830\to
830|943|830|109|252|252\to
903|850|850|252|109|252.
\end{align*}

\subsection{Closure of \texorpdfstring{\(U_\pi^{(3)}\)}{U-pi 3}}
\begin{align*}
&903|850|850|252|109|252\to
292|189|232|903|850|850\to\\
&292|189|232|923|810|870\to
870|923|810|169|272|272\to\\
&272|169|272|963|890|850\to
850|963|890|149|252|212\to\\
&212|149|252|943|830|830\to
830|943|830|109|252|252\to
903|850|850|252|109|252.
\end{align*}

\subsection{Closure of \texorpdfstring{\(U_\pi^{(4)}\)}{U-pi 4}}
\begin{align*}
&903|850|850|252|109|252\to
272|169|272|903|850|850\to\\
&272|169|272|963|890|850\to
850|963|890|149|252|212\to\\
&212|149|252|923|810|870\to
870|923|810|418|674|521\to\\
&943|830|830|521|418|674\to
830|943|830|109|252|252\to
903|850|850|252|109|252.
\end{align*}

\subsection{Closure of \texorpdfstring{\(U_\pi^{(5)}\)}{U-pi 5}}
\begin{align*}
&252|109|252|903|850|850\to
903|850|850|252|109|252\to\\
&614|498|501|252|109|252\to
501|614|498|169|272|272\to\\
&272|169|272|963|890|850\to
850|963|890|149|252|212\to\\
&212|149|252|943|830|830\to
830|943|830|109|252|252\to
252|109|252|903|850|850.
\end{align*}

\begin{theorem}[Minimality of the Five Closures]
\label{gfp:b:thm:minimalidad-cierres-pi-autonoma}
For each \(i\in\{1,\ldots,5\}\), the minimum length of a closed walk
that simultaneously traverses
\[
 U_\pi^{(i)},\quad U_e,\quad U_\varphi,\quad U_{\rm APP}
\]
in the quotient multigraph is eight.
\end{theorem}

\begin{proof}
Given \(i\), let
\[
 E_*=\{U_\pi^{(i)},U_e,U_\varphi,U_{\rm APP}\}.
\]
The finite space is considered.
\[
 \mathscr S_i=\{(v,M):v\in V(G),\ M\subseteq E_*\}.
\]
The first component records the vertex and the second the target edges
already traversed. If an edge \(e\) takes \(v\) to \(v'\), the transition is
\[
 (v,M)\longmapsto
 \bigl(v',M\cup(E_*\cap\{e\})\bigr).
\]
Breadth-first search enumerates lengths in increasing order. The first
return to \((v,E_*)\) occurs at length eight for all five regions.
The printed walks attain the bound; the optimality of the search excludes
a shorter length.
\end{proof}

The closure function is encoded by
\[
 \mathfrak C_\pi=
 \bigl(\mathscr R_\pi^{\rm micro},
 w_\pi,\rho_\pi,m_\pi,\ell_{\rm cl}\bigr),
\]
\[
 \rho_\pi=(\perp,++ ,++ ,++ ,--),
 \qquad
 m_\pi=(432,144,144,144,144),
 \qquad
 \ell_{\rm cl}=8.
\]

\section{Incidence operator and oriented half-turn}

The incidence shadow is
\[
 A_W=
 \begin{pmatrix}
 0&1&1&1&1&1\\
 1&0&1&2&2&1\\
 1&1&0&1&2&2\\
 1&2&1&0&1&2\\
 1&2&2&1&0&1\\
 1&1&2&2&1&0
 \end{pmatrix}.
\]
Multiplication in \(\mathbb F_3\) gives
\begin{equation}
 \boxed{A_W^2=2I_6=-I_6.}
 \label{gfp:b:eq:aw-cuadrado-menos-identidad-pi}
\end{equation}
Products between distinct rows and columns vanish modulo three,
whereas each diagonal product equals two. One application effects a ternary
quarter-turn; two applications effect the oriented inversion.

\section{Pending \texorpdfstring{\(1/5\)}{1/5} determined by the regional partition \(432+4\cdot144=1008\)}

Let
\[
 \lambda=(\lambda_1,\ldots,\lambda_5)^{\mathsf T}
 \in\mathbb Q_{\geq0}^5,
 \qquad
 \mathbf1^{\mathsf T}\lambda=1.
\]
Forgetting the regional label applies the symmetrizer
\[
 P_5=\frac15\mathbf1\mathbf1^{\mathsf T}.
\]
For every normalized vector,
\[
 P_5\lambda=\frac15\mathbf1.
\]
The invariance \(P_5\lambda=\lambda\) therefore has a unique solution:
\begin{equation}
 \boxed{\lambda=\frac15\mathbf1,\qquad t_1=\frac15.}
 \label{gfp:b:eq:pendiente-primitiva-pi}
\end{equation}
The integer five comes from the regional microfiber; it is not recognized in a
formula for \(\pi\).

\section{Rational composition and compensator \texorpdfstring{\(239\)}{239}}

The law of slope composition is
\[
 t\star s:=\frac{t+s}{1-ts}.
\]
The first duplication gives
\[
 t_2=t_1\star t_1
 =\frac{2/5}{1-1/25}
 =\frac5{12}.
\]
The second gives
\[
 t_4=t_2\star t_2
 =\frac{10/12}{1-25/144}
 =\frac{120}{119}.
\]
The return to the unit diagonal requires \(q>0\) such that
\[
 \frac{120/119-1/q}{1+(120/119)/q}=1.
\]
Multiplying by \(119q\) produces
\[
 120q-119=119q+120,
\]
and forces
\begin{equation}
 \boxed{q=239.}
 \label{gfp:b:eq:compensador-239-pi}
\end{equation}

\section{Gaussian identity and \(1^{\mathrm a}\) oriented half-turn}

We define
\begin{equation}
 \Theta_\pi
 :=16\arctan\frac15-4\arctan\frac1{239}.
 \label{gfp:b:eq:theta-pi-autonoma}
\end{equation}
The successive squares are
\[
\begin{array}{c|r@{}l}
\text{power}&\multicolumn{2}{c}{\text{Gaussian integer}}\\ \hline
(5+i)^2&24&+10i\\
(5+i)^4&476&+480i\\
(5+i)^8&-3824&+456960i\\
(5+i)^{16}&-208797818624&-3494830080i\\
(239-i)^2&57120&-478i\\
(239-i)^4&3262465916&-54606720i
\end{array}
\]
and the final multiplication gives
\begin{equation}
 \boxed{
 (5+i)^{16}(239-i)^4
 =-681386607803576157184.}
 \label{gfp:b:eq:identidad-gaussiana-pi-autonoma}
\end{equation}
The imaginary part is zero and the real part is negative. Therefore,
\[
 \Theta_\pi\equiv\pi\pmod{2\pi}.
\]
Para \(0<x<1\),
\[
 x-\frac{x^3}{3}<\arctan x<x.
\]
From here
\[
 \Theta_\pi>
 16\left(\frac15-\frac1{3\cdot5^3}\right)-\frac4{239}>3
\]
and
\[
 \Theta_\pi<\frac{16}{5}<4.
\]
The only negative orientation in that interval is the first positive half-turn:
\begin{equation}
 \boxed{\Theta_\pi=\pi.}
 \label{gfp:b:eq:reconocimiento-pi-exacto}
\end{equation}

The order of subsequent Archimedean recognition is
\[
 \text{five regions}
 \to\frac15
 \to\frac5{12}
 \to\frac{120}{119}
 \to239
 \to\Theta_\pi
 \to\pi.
\]

\section{Arbitrary precision rational horocycles}

For \(0<x<1\), let
\[
 S_n(x)=\sum_{r=0}^{n}(-1)^r\frac{x^{2r+1}}{2r+1}.
\]
Leibniz's criterion gives
\[
 S_{2m+1}(x)<\arctan x<S_{2m}(x).
\]
We define
\[
 a_m^-=S_{2m+1}(1/5),\quad a_m^+=S_{2m}(1/5),
\]
\[
 b_m^-=S_{2m+1}(1/239),\quad b_m^+=S_{2m}(1/239),
\]
and
\begin{equation}
 \ell_{\pi,m}=16a_m^- -4b_m^+,
 \qquad
 h_{\pi,m}=16a_m^+ -4b_m^-.
 \label{gfp:b:eq:horquilla-pi-autonoma}
\end{equation}
Then
\[
 \ell_{\pi,m}<\pi<h_{\pi,m}
\]
and
\begin{equation}
 h_{\pi,m}-\ell_{\pi,m}
 =
 \frac{16}{(4m+3)5^{4m+3}}
 +\frac{4}{(4m+3)239^{4m+3}}
 \longrightarrow0.
 \label{gfp:b:eq:anchura-horquilla-pi}
\end{equation}
The lower bounds grow and the upper bounds decrease.

Since \(3<\pi<4\), the ternary reader acts on
\[
 r_\pi:=\pi-3\in(0,1),
\]
with bracket
\[
 [\ell_{\pi,m}-3,h_{\pi,m}-3].
\]

\section{Subsequent Archimedean certification of the already generated fiber}

Suppose the internal relation \(\operatorname{Ext}\) and coinductive
survival have already produced, without using Machin or a tabulated expansion,
the compatible extension \(N_{K+1}\) of the prefix \(N_K\), of length
\(L_K=6K\). The \(729\) subcylinders of its ambient fiber are
\[
 I_{K,u}=
 \left[
 \frac{729N_K+u}{3^{L_K+6}},
 \frac{729N_K+u+1}{3^{L_K+6}}
 \right),
 \qquad0\leq u<729.
\]
The order \(m=m(K)\) is then increased until the Leibniz bracket
locates the already produced subcylinder. Let
\[
 \ell_{\pi,K}:=\ell_{\pi,m(K)}-3,
 \qquad
 h_{\pi,K}:=h_{\pi,m(K)}-3.
\]
The compatible integer endpoints are
\[
 j_{\min}=
 \max\!\left(
 729N_K,\left\lfloor
 \ell_{\pi,K}3^{L_K+6}\right\rfloor
 \right),
\]
\[
 j_{\max}=
 \min\!\left(
 729N_K+728,\left\lceil
 h_{\pi,K}3^{L_K+6}\right\rceil-1
 \right).
\]
The condition of unitary localization is
\begin{equation}
 \boxed{j_{\min}=j_{\max},\qquad N_{K+1}=j_{\min}.}
 \label{gfp:b:eq:seleccion-unitaria-subcilindro-pi}
\end{equation}
The first equality certifies unit localization; the second verifies
that this localization agrees with the internally generated extension. Neither
defines the TPK transition or selects among the \(729\) states.
Existence and uniqueness come from \(\operatorname{Ext}\) and
survival; Machin and the Leibniz brackets are subsequent
recognitions.

\section{Coinductive prolongation of the section of \(\pi\) through the inverse system of surviving histories}

Let \(\mathcal H_m^\pi\) be the set of enriched histories at depth
\(m\), and let
\[
 p_{n,m}:\mathcal H_n^\pi\longrightarrow\mathcal H_m^\pi
\]
the truncation. We define
\[
 \operatorname{Surv}_m^{(N),\pi}
 :=p_{N,m}(\mathcal H_N^\pi),
\]
\[
 \operatorname{Surv}_m^\pi
 :=\bigcap_{N\geq m}\operatorname{Surv}_m^{(N),\pi}.
\]
When the relation \(\operatorname{Ext}\) is nonempty, single-valued and
compatible at every level, the sets are singletons and satisfy
\[
 p_{m+1,m}(\operatorname{Surv}_{m+1}^{\pi})
 =\operatorname{Surv}_{m}^{\pi}.
\]
In that domain of extendability, there exists, therefore, a unique section.
\[
 x_\pi\in
 X_\pi:=\varprojlim_m\operatorname{Surv}_m^\pi.
\]
The cylinders are nested and their diameters are \(3^{-6K}\), which tends to
zero. Their intersection contains a unique point, identified by the angular
character with \(r_\pi=\pi-3\).

For every decimal depth \(M\), one may choose a level \(K(M)\) whose
cylinder lies within a single decimal cell of diameter \(10^{-M}\).
The deep prefixes truncate exactly to the preceding ones.

\section{Structural characterization, necessary validity conditions and obstructions}

\begin{theorem}[Characterization of the closure channel]
\label{gfp:b:thm:caracterizacion-final-pi-autonoma}
Starting from the catalog, the action \(D_3\), the oriented gauge, the
operators \(L_0,L_1,A_W\), the slope law and the rational filtration:
\begin{enumerate}
 \item \(w_6^\pi=010211\) is selected;
 \item five regions are constructed with the identities of
       \eqref{gfp:b:eq:descomposiciones-cinco-regiones-pi};
 \item each region belongs to a minimum closure of length eight;
 \item \(A_W^2=-I_6\) fixes the oriented inversion;
 \item regional symmetrization forces \(1/5\);
 \item the composition forces \(120/119\);
 \item the unit return forces \(239\);
 \item the Gaussian identity recognizes the first half-turn;
 \item the brackets and the fiber of \(729\) elements produce a section
       compatible at every finite depth.
\end{enumerate}
Its evaluation is
\[
 \boxed{\pi=16\arctan\frac15-4\arctan\frac1{239}.}
\]
\end{theorem}

The complete chain is
\[
 \mathrm{APP}\to\mathrm{TRIT}\to\mathrm{TPK}
 \to104\,976\to468\to243\to w_6^\pi
 \to w_{30}^\pi\to R_{36}\to\Gamma_9
 \to X_\pi\to\pi.
\]

\begin{remark}[Certified and validity control: Forgery of the closure channel]
The derivation is refuted if:
\begin{enumerate}
 \item the closure orbit ceases to be unique;
 \item one of the five regions disappears or its complete identity changes;
 \item one of the five closures has length less than eight;
 \item \(A_W^2\ne2I_6\);
 \item the symmetrizer admits another normalized distribution;
 \item the compositions do not produce \(5/12\), \(120/119\), and \(239\);
 \item the Gaussian identity fails;
 \item a bracket ceases to contain the angular character;
 \item a level does not traverse the \(729\) elements of the ambient fiber or does not
       locate exactly one extension;
 \item a target expansion intervenes before recognition.
\end{enumerate}
\end{remark}
